\chapter{Conclusions and Future Work}
\label{chap:conclusions}

This chapter summarises the key achievements of the AI-Assisted ETF Recommendation Platform, discusses the limitations encountered during development, and outlines potential directions for future enhancement.

%-------------------------------------------------------------------
\section{Summary of Achievements}
\label{sec:achievements}

This project successfully delivered a comprehensive web-based platform that empowers retail investors to navigate the complex ETF landscape with confidence. The key accomplishments are summarised below.

\textbf{Technical Achievements:}
\begin{itemize}
    \item Developed a full-stack web application using Next.js 14, FastAPI, and PostgreSQL, following modern software engineering practices.
    \item Aggregated a comprehensive ETF database covering 2,272 ETFs from the US and Singapore markets, with over 5.7 million historical price records spanning 15 years.
    \item Implemented a real-time metrics calculation engine computing beta, Sharpe ratio, volatility, and maximum drawdown on-demand using vectorised Pandas operations.
    \item Integrated an AI-powered chatbot leveraging OpenAI's GPT-4o model to provide plain-language explanations of ETF concepts with appropriate disclaimers.
    \item Created a multi-factor recommendation engine that scores ETFs based on risk alignment (40\%), performance (30\%), cost efficiency (15\%), and liquidity (15\%), with explainable score breakdowns.
    \item Developed scenario analysis tools featuring historical stress testing (COVID-19 crash replay) and Value at Risk (VaR) calculations at 95\% confidence level.
    \item Implemented real-time news aggregation with sentiment analysis via Alpha Vantage API integration.
\end{itemize}

\textbf{User Experience Achievements:}
\begin{itemize}
    \item Designed an intuitive, mobile-responsive interface accessible to users with varying levels of financial knowledge.
    \item Created an interactive 15-question risk profiling questionnaire that classifies users into four distinct risk profiles (Conservative, Moderate, Growth, Aggressive).
    \item Achieved a user satisfaction rating of 4.3/5.0 in User Acceptance Testing with 8 participants.
    \item Delivered professional-grade analytics typically reserved for institutional investors in an accessible format for retail users.
\end{itemize}

\textbf{Project Management:}
All six core objectives outlined in Chapter 1 were successfully achieved, with the platform deployed and fully operational on Vercel (frontend), Render (backend), and Supabase (database).

%-------------------------------------------------------------------
\section{Limitations}
\label{sec:limitations}

Despite the project's successes, several limitations were identified that constrain the platform's capabilities.

\textbf{Data Limitations:}
\begin{itemize}
    \item Market coverage is limited to US and Singapore ETFs; major markets such as Europe, Japan, and Hong Kong are not included.
    \item Price data is updated daily rather than in real-time, which may not reflect intraday price movements.
    \item Historical data gaps exist for some smaller or newer ETFs with limited trading history.
    \item The free-tier API rate limits (Yahoo Finance, Alpha Vantage) occasionally affect data freshness during peak usage.
\end{itemize}

\textbf{Algorithmic Limitations:}
\begin{itemize}
    \item The recommendation engine uses a rule-based multi-factor scoring approach rather than machine learning models, which may not capture complex non-linear relationships.
    \item VaR calculations assume normally distributed returns, which underestimates tail risk during extreme market events.
    \item Sentiment analysis accuracy is dependent on the Alpha Vantage API's underlying NLP models.
\end{itemize}

\textbf{Scope Limitations:}
\begin{itemize}
    \item The platform provides educational information only and does not constitute regulated financial advice.
    \item No trading execution capability is integrated; users must execute trades through external brokerages.
    \item Multi-currency support is not implemented, limiting usability for international investors.
\end{itemize}

%-------------------------------------------------------------------
\section{Future Work}
\label{sec:future-work}

Several enhancements are proposed to extend the platform's capabilities and address the identified limitations.

\textbf{Short-term Enhancements (3--6 months):}
\begin{itemize}
    \item \textbf{ML-based Recommendations:} Replace the rule-based scoring algorithm with a machine learning model (e.g., gradient boosting or neural collaborative filtering) trained on user interaction data.
    \item \textbf{Real-time Price Streaming:} Implement WebSocket connections for live price updates during market hours.
    \item \textbf{Email Alerts:} Enable users to set price movement thresholds and receive automated email notifications.
    \item \textbf{Enhanced Portfolio Analytics:} Add correlation matrix visualisation and Markowitz mean-variance optimisation suggestions.
\end{itemize}

\textbf{Medium-term Enhancements (6--12 months):}
\begin{itemize}
    \item \textbf{Historical Backtesting:} Allow users to simulate portfolio performance against historical data with custom start dates and rebalancing strategies.
    \item \textbf{Additional Markets:} Expand coverage to include European (UCITS) and Asia-Pacific ETFs.
    \item \textbf{Multi-currency Support:} Display prices and returns in the user's preferred currency with automatic conversion.
    \item \textbf{Dividend Analysis:} Integrate dividend history, yield calculations, and reinvestment projections.
\end{itemize}

\textbf{Long-term Vision (1--2 years):}
\begin{itemize}
    \item \textbf{Mobile Application:} Develop native iOS and Android applications using React Native for improved mobile experience.
    \item \textbf{Social Features:} Enable users to share portfolios, follow other investors, and participate in community discussions.
    \item \textbf{Brokerage Integration:} Partner with discount brokerages to enable seamless trade execution directly from the platform.
    \item \textbf{Regulatory Exploration:} Investigate requirements for obtaining financial advisory licensing in target jurisdictions.
\end{itemize}

%-------------------------------------------------------------------
\section{Final Remarks}
\label{sec:final-remarks}

This project demonstrates that modern web technologies combined with artificial intelligence can democratise access to sophisticated investment analysis tools. By providing retail investors with explainable recommendations, educational AI assistance, and professional-grade analytics, the platform bridges the gap between institutional-quality tools and everyday investors.

The development process reinforced the importance of user-centric design in financial applications, where clarity and trust are paramount. The integration of appropriate disclaimers ensures that users understand the educational nature of the platform while still benefiting from its analytical capabilities.

As AI continues to transform the financial services industry, platforms like this have the potential to significantly improve financial literacy and empower individuals to make more informed investment decisions. The author hopes that this project contributes meaningfully to the ongoing democratisation of financial technology and inspires further innovation in accessible investment education.
